\documentclass[11pt]{article}

\usepackage[margin=1in]{geometry}
\usepackage{amsmath,amssymb,amsthm,mathtools}
\usepackage{hyperref}
\usepackage{microtype}

\newtheorem{theorem}{Theorem}[section]
\newtheorem{corollary}[theorem]{Corollary}
\newtheorem{proposition}[theorem]{Proposition}
\newtheorem{lemma}[theorem]{Lemma}
\theoremstyle{remark}
\newtheorem{remark}[theorem]{Remark}

\newcommand{\C}{\mathbb C}
\newcommand{\R}{\mathbb R}
\newcommand{\conv}{\operatorname{conv}}
\newcommand{\cl}{\operatorname{cl}}
\newcommand{\interior}{\operatorname{int}}

\title{Cofinal Convex Recurrence of the Riemann-Zero Wave Trace\\
and Suzuki's Screw Function}
\author{}
\date{Draft of July 22, 2026}

\begin{document}
\maketitle

\begin{abstract}
We prove an unconditional recurrence law for a hierarchy of regularized wave traces formed from
the nontrivial zeros of the Riemann zeta function.  If $\Gamma$ is the multiplicity-exact zero
frequency multiset of the centered xi-function, set
\[
  H_j(t)=\sum_{\gamma\in\Gamma}
       \frac{\sin(\gamma t)}{\gamma^{2j+1}},\qquad j\geq 1.
\]
For every $K\geq1$ and every $T>0$, the origin is an interior point of the closed convex hull of
$\{(H_1(t),\ldots,H_K(t)):t>T\}$.  Consequently, at most $K+1$ distinct times beyond $T$ carry
positive weights that annihilate all first $K$ traces simultaneously.  No Riemann hypothesis,
simplicity assumption, or rational independence of zero ordinates is used.

The proof is a universal pole-frontier principle.  For a symmetric discrete spectrum of finite
horizontal width, eventual one-sidedness of a polynomially weighted odd trace forces a purely
imaginary spectral point on its active rightmost pole frontier.  This obstruction is sharp.  The
centered xi-divisor has no such point because xi has no real nontrivial zero; every nonzero linear
combination of the finite trace vector therefore escapes both half-spaces cofinally.  Separation
and Carath\'eodory's theorem yield the positive quadrature.

The first trace is the curvature-centered primitive of Suzuki's screw function.  Thus the scalar
case gives an unconditional two-sided crossing of its initial C\`esaro mean, while the vector
theorem supplies a new all-order simultaneous recurrence law for the Suzuki--xi spectral
hierarchy.  The result is not a proof of the Riemann hypothesis; it gives source-side structure for
the finite-aperture spectral program and isolates real-edge divisors as the only universal
obstruction to this recurrence mechanism.
\end{abstract}

\section{Introduction}

Suzuki associated a continuous screw function with the completed Riemann zeta function and used
it to express Weil positivity, Li-type criteria, and canonical spectral constructions in a common
framework~\cite{Suzuki2023}.  A recent finite-interval formulation realizes restricted Weil forms
through nonlocal self-adjoint extensions of a first-order differential operator and proposes a
cofinal spectral limit toward the zero ordinates~\cite{Suzuki2026}.  These developments make the
time dependence of the screw profile, rather than only its distributional Fourier transform, a
natural object of study.

Let
\[
  \Xi(z)=\xi\!\left(\frac12+iz\right),
\]
and let $\Gamma$ denote its zeros, counted with analytic multiplicity.  The functional equation
and Schwarz symmetry make $\Gamma$ invariant under $\gamma\mapsto-\gamma$ and complex
conjugation.  Suzuki's multiplicity-exact expansion has the form
\[
  \Psi(t)=\sum_{\gamma\in\Gamma}
       \frac{1-\cos(\gamma t)}{\gamma^2}.
  \tag{1.1}
\]
Under the Riemann hypothesis every frequency is real, and ordinary almost-periodic intuition
suggests recurrent oscillation.  Without the Riemann hypothesis, however, a zero away from the
critical line contributes an exponentially growing oscillatory mode.  A finite exponential sum
with a unique extremal mode can become one-signed, so unconditional recurrence is not a formal
consequence of the functional equation.

Our starting point is the odd inverse-power hierarchy
\[
  H_j(t)=\sum_{\gamma\in\Gamma}
       \frac{\sin(\gamma t)}{\gamma^{2j+1}},
  \qquad j=1,2,\ldots.
  \tag{1.2}
\]
These functions are regularized wave traces.  Their static inverse moments belong to the
superzeta framework developed by Voros~\cite{Voros2003,Voros2014}, while their time dependence is
adapted to the screw function.  The hierarchy obeys
\[
  H_{j+1}''=-H_j,\qquad
  H_j(0)=0,\qquad
  H_j'(0)=\sum_{\gamma\in\Gamma}\gamma^{-2j}.
  \tag{1.3}
\]

The main result is a finite-dimensional recurrence theorem stronger than coordinatewise sign
change.  For every finite initial segment of (1.2), every nonzero real linear functional assumes
both strict signs after every time.  Equivalently, the trace vector eventually lies on both sides
of every hyperplane through the origin.  This half-space escape implies that the origin is an
interior point of every cofinal closed convex hull.  A finite-dimensional convexity argument then
produces an exact positive atomic quadrature supported at no more than $K+1$ arbitrarily late
times.

The mechanism is universal and sharp.  A polynomial linear combination of (1.2) has a Laplace
transform whose genuine poles are the active spectral frequencies.  If the time function were
eventually one-signed, Landau's boundary theorem would force a singularity at the real point of
its convergence abscissa.  The meromorphic pole expansion is regular at that point unless an
active frequency is purely imaginary and lies at the spectral edge.  A pair of frequencies
$\{i\delta,-i\delta\}$ shows that this exception cannot be removed: it produces a hyperbolic-sine
mode of fixed sign.

For the Riemann xi-function, a purely imaginary $\gamma$ would correspond to a real nontrivial
zero of xi in $(0,1)$, which does not exist.  This elementary real-axis fact, not the Riemann
hypothesis, removes the sole universal obstruction.  The result therefore remains valid even if
there are zeros away from the critical line.

Section~\ref{sec:universal} states and proves the pole-frontier theorem.  Section~\ref{sec:convex}
derives cofinal convex recurrence and positive quadrature.  Section~\ref{sec:xi} verifies the xi
hypotheses and identifies $H_1$ with Suzuki's curvature-centered screw primitive.  The final
section records the contribution and its limits for the finite-aperture RH program.

\section{A universal pole-frontier recurrence theorem}
\label{sec:universal}

We work with multisets; every sum includes multiplicity.

\subsection{Admissible symmetric spectra}

Let $\Gamma$ be an infinite locally finite multiset in $\C\setminus\{0\}$.  We call $\Gamma$
\emph{admissible} if:
\begin{enumerate}
\item $\Gamma$ is invariant, with multiplicity, under $\gamma\mapsto-\gamma$ and
      $\gamma\mapsto\overline\gamma$;
\item $\sum_{\gamma\in\Gamma}|\gamma|^{-2}<\infty$;
\item $\delta:=\sup_{\gamma\in\Gamma}|\operatorname{Im}\gamma|<\infty$;
\item the set $\{\gamma^2:\gamma\in\Gamma\}$ is infinite.
\end{enumerate}

For a real polynomial $P$, define its active multiset and active width by
\[
 \Gamma_P=\{\gamma\in\Gamma:P(\gamma^{-2})\neq0\},
 \qquad
 \delta_P=\sup_{\gamma\in\Gamma_P}|\operatorname{Im}\gamma|,
 \tag{2.1}
\]
and define
\[
  H_P(t)=\sum_{\gamma\in\Gamma}
  \frac{\sin(\gamma t)}{\gamma^3}P(\gamma^{-2}).
  \tag{2.2}
\]
Because a nonzero polynomial has finitely many roots while the squared spectrum is infinite,
$\Gamma_P$ is infinite whenever $P\neq0$.

\begin{lemma}[Normal convergence]
\label{lem:normal}
For every real polynomial $P$, the series in \emph{(2.2)} converges normally on compact subsets of
$\R$ and defines a real continuous function satisfying
\[
 |H_P(t)|\leq C_P e^{\delta_P t},\qquad t\geq0.
 \tag{2.3}
\]
\end{lemma}

\begin{proof}
Local finiteness and $0\notin\Gamma$ give a positive lower bound for $|\gamma|$.  Hence reciprocal
square summability implies reciprocal cube summability.  The values $P(\gamma^{-2})$ are bounded,
and
\[
  |\sin(\gamma t)|\leq e^{|\operatorname{Im}\gamma|t}.
\]
The Weierstrass test proves normal convergence and (2.3).  Conjugation invariance pairs every
summand with its complex conjugate, proving reality.
\end{proof}

\begin{lemma}[Active-pole transform]
\label{lem:transform}
For $\operatorname{Re}s>\delta_P$,
\[
 \mathcal L H_P(s)
 =M_P(s):=
 \sum_{\gamma\in\Gamma}
 \frac{P(\gamma^{-2})}{\gamma^2(s^2+\gamma^2)}.
 \tag{2.4}
\]
The right-hand side extends meromorphically to $\C$.  Its poles are exactly
$s=\pm i\gamma$ for $\gamma\in\Gamma_P$, with the identifications forced by
$\gamma\leftrightarrow-\gamma$, and every such pole is genuine.
\end{lemma}

\begin{proof}
For $\operatorname{Re}s>\delta_P$, Lemma~\ref{lem:normal} permits termwise integration and
\[
 \int_0^\infty e^{-st}\sin(\gamma t)\,dt
 =\frac{\gamma}{s^2+\gamma^2}.
\]
This gives (2.4).  On a compact set avoiding the displayed poles, every sufficiently large
summand is $O_P(|\gamma|^{-4})$; normal convergence therefore gives a meromorphic function on
$\C$.  The $\gamma$ and $-\gamma$ terms are identical, so their residues add.  No other frequency
has the same square, and an active residue cannot vanish.
\end{proof}

\begin{theorem}[Pole-frontier recurrence dichotomy]
\label{thm:frontier}
Let $\Gamma$ be admissible and let $P\in\R[x]$ be nonzero.  If $H_P$ is eventually nonnegative or
eventually nonpositive, then the active frontier contains a purely imaginary frequency:
\[
  i\delta_P\in\Gamma_P
  \quad\text{or}\quad
  -i\delta_P\in\Gamma_P.
  \tag{2.5}
\]
Consequently, if $\Gamma$ has no purely imaginary frequency, then for every $T>0$ there exist
$t_+,t_->T$ such that
\[
 H_P(t_-)<0<H_P(t_+).
 \tag{2.6}
\]
\end{theorem}

\begin{proof}
Suppose first that $H_P$ is eventually nonnegative.  Choose $T$ such that
$h_T(u):=H_P(T+u)\geq0$ for $u\geq0$.  Its Laplace transform agrees initially with
\[
 M_{P,T}(s)=e^{sT}\left(
 M_P(s)-\int_0^T H_P(u)e^{-su}\,du\right).
 \tag{2.7}
\]
The finite integral is entire, so Lemma~\ref{lem:transform} shows that (2.7) has precisely the
active poles of $M_P$.

Let $\sigma_c$ be the convergence abscissa of the nonnegative Laplace integral of $h_T$.  Bound
(2.3) gives $\sigma_c\leq\delta_P$.  If $\sigma_c<\delta_P$, the definition of $\delta_P$
provides an active $\gamma$ for which one of $i\gamma$ and $-i\gamma$ has real part greater than
$\sigma_c$.  The Laplace integral is holomorphic at that point, whereas its meromorphic
continuation (2.7) has a genuine pole there.  This is impossible.  Hence
\[
  \sigma_c=\delta_P.
  \tag{2.8}
\]
The tail $h_T$ is not identically zero, since otherwise (2.7) would have no pole.

Landau's boundary theorem for Laplace transforms of nonzero nonnegative functions forces a
singularity at the real convergence-boundary point $s=\sigma_c$; see, for example,
Widder~\cite[Chapter II, Theorem 5b]{Widder1941}.  Local finiteness of $\Gamma$ and the normal
convergence in Lemma~\ref{lem:transform} show that $M_P$ is holomorphic in a neighborhood of the
real point $\delta_P$ unless this point itself is an active pole.  In view of (2.8), the required
singularity therefore implies
\[
 \delta_P=\pm i\gamma
\]
for an active $\gamma$, which is exactly (2.5).  The entire correction in (2.7) cannot alter this
conclusion.

Applying the same argument to $-H_P$ treats eventual nonpositivity.  If no frequency is purely
imaginary, neither eventual sign is possible.  Continuity then gives (2.6) after every $T$.
\end{proof}

\begin{remark}[Sharpness]
The exceptional alternative in Theorem~\ref{thm:frontier} cannot be removed.  If
$\Gamma=\{i\delta,-i\delta\}$ and $P=1$, counting the symmetric pair once each, then
\[
 H_P(t)=\frac{2(-1)^1}{\delta^3}\sinh(\delta t)
\]
up to the harmless parity sign dictated by the chosen inverse power.  It is one-signed for
$t>0$, and the Laplace transform has its required real boundary pole at $s=\delta$.
\end{remark}

\section{Cofinal convex recurrence and positive quadrature}
\label{sec:convex}

Assume throughout this section that $\Gamma$ is admissible and contains no purely imaginary
frequency.  For $j\geq1$, set
\[
 H_j(t)=\sum_{\gamma\in\Gamma}
       \frac{\sin(\gamma t)}{\gamma^{2j+1}},
 \qquad
 V_K(t)=(H_1(t),\ldots,H_K(t))\in\R^K.
 \tag{3.1}
\]

\begin{theorem}[Cofinal convex recurrence]
\label{thm:convex}
For every $K\geq1$ and every $T>0$,
\[
 0\in\interior\cl\conv\{V_K(t):t>T\}.
 \tag{3.2}
\]
Equivalently, the cofinal trace vector is not contained in any closed half-space whose boundary
passes through the origin.
\end{theorem}

\begin{proof}
For $c=(c_1,\ldots,c_K)\in\R^K$, define
\[
 P_c(x)=\sum_{j=1}^K c_jx^{j-1}.
 \tag{3.3}
\]
Then
\[
 c\cdot V_K(t)=H_{P_c}(t).
 \tag{3.4}
\]
If $c\neq0$, then $P_c\neq0$, and Theorem~\ref{thm:frontier} says that (3.4) assumes both strict
signs for $t>T$.

Let $C_T$ denote the closed convex set in (3.2).  If the origin were not an interior point of
$C_T$, finite-dimensional separation would produce a nonzero $c\in\R^K$ such that, after replacing
$c$ by $-c$ if necessary,
\[
 c\cdot x\geq0\qquad(x\in C_T).
\]
This contradicts the negative value of (3.4) at some $t>T$.  Therefore (3.2) holds.
\end{proof}

\begin{corollary}[Positive cofinal annihilating quadrature]
\label{cor:quadrature}
For every $K\geq1$ and $T>0$, there exist an integer $m$ with $1\leq m\leq K+1$, distinct times
$t_1,\ldots,t_m>T$, and weights $a_1,\ldots,a_m>0$ such that
\[
 \sum_{r=1}^m a_r=1,
 \qquad
 \sum_{r=1}^m a_rH_j(t_r)=0
 \quad(1\leq j\leq K).
 \tag{3.5}
\]
\end{corollary}

\begin{proof}
Taking closure does not change the relative interior of a convex subset of a finite-dimensional
real vector space~\cite[Section 6]{Rockafellar1970}.  Theorem~\ref{thm:convex} gives full interior,
so the origin already belongs to the ordinary convex hull of the tail.  Carath\'eodory's theorem
represents it using at most $K+1$ points.  Delete zero weights and merge repeated times to obtain
(3.5).
\end{proof}

The conclusion is genuinely simultaneous.  Separate sign changes of the coordinates $H_j$ would
not imply (3.2) or (3.5); the polynomial identity (3.4) is the step that controls every direction
in the finite trace space.

\section{The centered xi-divisor and Suzuki's screw function}
\label{sec:xi}

We now take
\[
 \Xi(z)=\xi\!\left(\frac12+iz\right)
\]
and let $\Gamma$ be its zero multiset.  The standard properties required below may be found in
classical accounts of the zeta function, for example~\cite{Titchmarsh1986}.

\begin{proposition}[Admissibility of the xi frequency spectrum]
\label{prop:xi-admissible}
The multiset $\Gamma$ is admissible and contains no purely imaginary frequency.
\end{proposition}

\begin{proof}
The functional equation makes $\Xi$ even, while Schwarz symmetry makes its divisor invariant under
conjugation.  The centered xi-function is entire of order one; its zeros are locally finite,
infinite, and satisfy reciprocal-square summability.  The critical strip gives
\[
 |\operatorname{Im}\gamma|<\frac12.
\]
Infinitely many nontrivial zeros with unbounded ordinates imply infinitely many distinct values of
$\gamma^2$.

If $\gamma=iy$ were purely imaginary, then
\[
 \Xi(iy)=\xi\!\left(\frac12-y\right)=0
\]
would be a real nontrivial zero in $(0,1)$.  The zeta function has no real zero in that interval,
so no such $\gamma$ exists.  Also $\gamma=0$ is excluded because $\xi(1/2)\neq0$.
\end{proof}

Combining Proposition~\ref{prop:xi-admissible} with Theorem~\ref{thm:convex} and
Corollary~\ref{cor:quadrature} proves the principal arithmetic statement.

\begin{theorem}[Xi zero-trace recurrence]
\label{thm:xi}
For every $K\geq1$ and $T>0$, the vector of xi zero traces in \emph{(3.1)} satisfies \emph{(3.2)}
and admits the positive annihilating quadrature \emph{(3.5)}.  The assertion is unconditional.
\end{theorem}

To connect the first coordinate with Suzuki's screw function, use the convention
$g=-\Psi$ and write
\[
 c_\Xi=\frac{\Xi''(0)}{\Xi(0)}
      =-\sum_{\gamma\in\Gamma}\frac1{\gamma^2}.
 \tag{4.1}
\]
Suzuki's expansion (1.1) gives
\[
 g(t)-c_\Xi
 =\sum_{\gamma\in\Gamma}\frac{\cos(\gamma t)}{\gamma^2}.
 \tag{4.2}
\]
Normal convergence permits termwise integration:
\[
 \int_0^A(g(t)-c_\Xi)\,dt
 =\sum_{\gamma\in\Gamma}\frac{\sin(\gamma A)}{\gamma^3}
 =H_1(A).
 \tag{4.3}
\]

\begin{corollary}[Unconditional screw-C\`esaro crossing]
\label{cor:cesaro}
For every $A_0>0$, there exist $A_-,A_+>A_0$ such that
\[
 \frac1{A_-}\int_0^{A_-}g(t)\,dt<c_\Xi
 <\frac1{A_+}\int_0^{A_+}g(t)\,dt.
 \tag{4.4}
\]
\end{corollary}

\begin{proof}
Apply Theorem~\ref{thm:frontier} to $P=1$, use (4.3), and divide the two strict signs by the
positive apertures.
\end{proof}

The higher coordinates are canonical rather than ad hoc.  Differentiating their normally
convergent series gives
\[
 H_{j+1}''=-H_j,qquad H_j(0)=0,qquad
 H_j'(0)=\sum_{\gamma\in\Gamma}\gamma^{-2j}.
 \tag{4.5}
\]
Thus Corollary~\ref{cor:quadrature} constructs, beyond every aperture threshold, a positive atomic
law supported on at most $K+1$ apertures that annihilates the first $K$ regularized primitives of
the centered screw curvature simultaneously.

For comparison with logarithmic-derivative and superzeta formulas, put
\[
 F(s)=\xi\!\left(\frac12+s\right),
 \qquad S_j=\sum_{\gamma\in\Gamma}\gamma^{-2j},
\]
and define
\[
 M_j(s)=\sum_{\gamma\in\Gamma}
 \frac1{\gamma^{2j}(s^2+\gamma^2)}.
\]
Then
\[
 M_0(s)=\frac{F'(s)}{sF(s)},
 \qquad
 M_j(s)=\frac{S_j-M_{j-1}(s)}{s^2}.
 \tag{4.6}
\]
Iterating (4.6) expresses every Laplace transform $\mathcal L H_j=M_j$ as a finite combination of
the inverse moments $S_1,\ldots,S_j$ and the centered logarithmic derivative $F'/F$.  The apparent
singularities at $s=0$ cancel by the even Taylor expansion of $F$.  Formula (4.6) places the dynamic
trace hierarchy next to the static superzeta hierarchy without assuming that the zeros are real.

\section{Significance and limits for the Suzuki RH lane}

Theorem~\ref{thm:xi} adds three pieces of unconditional structure to the screw-function program.

First, it identifies a universal obstruction.  Eventual one-sidedness of a regularized odd trace
cannot be caused merely by zeros off the critical line.  It requires a real-edge divisor, encoded
as a purely imaginary frequency at the active pole frontier.  The xi-function has no such divisor.

Second, the result is stable under every finite polynomial spectral filter.  A filter may delete
finitely many zero orbits and change the active horizontal width, but the Landau argument restarts
at the new frontier.  This stability is what upgrades the scalar C\`esaro crossing to finite-vector
half-space escape.

Third, the convex formulation produces exact positive cancellation.  For each finite regularized
jet order, a probability measure on at most $K+1$ arbitrarily large apertures annihilates all
corresponding traces.  Suzuki's finite-interval operators are indexed by aperture, so this gives a
canonical cofinal ensemble principle unavailable from a single sign-changing coordinate.

The theorem does not prove the Riemann hypothesis.  It neither proves positivity of Suzuki's full
restricted Weil kernel nor identifies a limiting self-adjoint operator with the zero divisor.  In
particular, positive cancellation across several apertures must not be confused with positivity at
one aperture.  Its role is source-side: it supplies an unconditional recurrence and quadrature law
that every proposed cofinal finite-aperture realization must respect.  A further theorem connecting
the positive aperture ensemble to the normalized prime--Gamma Pick kernel would be needed for an RH
argument.

\section*{Acknowledgment of prior structures}

The proof uses Suzuki's screw-function expansion~\cite{Suzuki2023}, the classical Landau boundary
theorem in the Laplace-transform form recorded by Widder~\cite{Widder1941}, standard finite-dimensional
convex separation, and Carath\'eodory's theorem.  Voros's superzeta theory~\cite{Voros2003,Voros2014}
provides the natural context for the inverse zero moments.  The new claim is the polynomially stable
pole-frontier dichotomy, its cofinal convex-hull consequence, and the exact positive simultaneous
quadrature for the Suzuki--xi trace hierarchy.  Priority for this combination should be checked by
independent specialist review.

\begin{thebibliography}{99}

\bibitem{Rockafellar1970}
R.~T. Rockafellar,
\emph{Convex Analysis},
Princeton Mathematical Series 28, Princeton University Press, 1970.

\bibitem{Suzuki2023}
M. Suzuki,
Aspects of the screw function corresponding to the Riemann zeta-function,
\emph{arXiv:2206.03682v4} (2023).

\bibitem{Suzuki2026}
M. Suzuki,
Weil's quadratic form via the screw function,
\emph{arXiv:2606.09096v1} (2026).

\bibitem{Titchmarsh1986}
E.~C. Titchmarsh,
\emph{The Theory of the Riemann Zeta-Function},
2nd ed., revised by D.~R. Heath-Brown, Clarendon Press, 1986.

\bibitem{Voros2003}
A. Voros,
Zeta functions over zeros of general zeta and $L$-functions,
\emph{arXiv:math/0311029} (2003).

\bibitem{Voros2014}
A. Voros,
Zeta functions over zeros of zeta functions and an exponential-asymptotic view of the Riemann
hypothesis,
\emph{RIMS K\^oky\^uroku Bessatsu} \textbf{B52} (2014), 147--164;
\emph{arXiv:1403.4558}.

\bibitem{Widder1941}
D.~V. Widder,
\emph{The Laplace Transform},
Princeton Mathematical Series 6, Princeton University Press, 1941.

\end{thebibliography}

\end{document}
